% Two-run comparisons reconstructed from the 16 validated starting-backlog trials.
% Individual observations and analysis provenance are retained with the experiment records.
The following figures compare the metrics in Table~\ref{tab:metric-roles}
using equal-weight means of the 2 runs for each response and backlog
condition. Bands and vertical whiskers show the range of those runs,
not confidence intervals. Time-series means require valid values from
both runs, so unavailable observations remain gaps. Gray dotted lines
mark the action request at minute~1, while colored dotted lines mark
confirmed activation in each run. The dashed throughput line marks
the continuing input rate of 600~messages/s.

Figure~\ref{fig:backlog-mean-control} shows that the larger-backlog
preparation also produced older waiting work. At the action request,
mean message age was approximately 42--43~seconds, compared with
14--15~seconds in the smaller-backlog condition. These ages include
acknowledged warm-up messages still awaiting completion, and they
cannot be separated from the effect of backlog size in this comparison.
Without intervention, backlog and message age continued increasing
in both conditions.

\begin{figure}[htbp]
\centering
\includegraphics[width=\linewidth,height=.78\textheight,keepaspectratio]{figures/backlog-mean-control.pdf}
\caption{Runtime control signals for smaller and larger starting backlogs.}
\label{fig:backlog-mean-control}
\end{figure}

The measured partition arrival rates remained similar across the two
backlog conditions, as shown in Figure~\ref{fig:backlog-mean-candidates}.
Before intervention, Consumer~0 received approximately
500~messages/s, exceeding its matching calibrated rate of about
287~messages/s. Redistribution's specified map spread this demand to approximately
200~messages/s per existing consumer. The capacity points are the fixed calibration
references used for these trials, rather than new estimates from each
repetition. Matching capacity measurements were unavailable for the
added consumers, so their capacities are left unreported.

\begin{figure}[htbp]
\centering
\includegraphics[width=\linewidth]{figures/backlog-mean-candidates.pdf}
\caption{Pre-action partition arrival rates and matching consumer capacities.}
\label{fig:backlog-mean-candidates}
\end{figure}

The throughput peaks in Figure~\ref{fig:backlog-mean-performance}
reflect processing of accumulated work while input continued.
Throughput includes distinct warm-up messages completed during
evaluation, whereas window p99 uses evaluation messages completed
within each 30-second interval. Figure~\ref{fig:backlog-mean-outcomes}
separately shows evaluation-message outcomes through the drain cutoff.
All intervention runs completed that cohort, while Keep~3 left a mean
of 31.0\% unfinished with the smaller backlog and 38.9\% with the larger
backlog. Run-level p99 points average the 2 run-level percentiles
without pooling message latencies.

\begin{figure}[htbp]
\centering
\includegraphics[width=\linewidth]{figures/backlog-mean-performance.pdf}
\caption{Throughput and window p99 during evaluation.}
\label{fig:backlog-mean-performance}
\end{figure}

Targeted scaling reduced mean recovery time from 239 to 174~seconds
for the smaller backlog and from 426 to 247~seconds for the larger
backlog, compared with redistribution. However, its mean run-level p99
was higher, increasing from 54.0 to 80.5~seconds and from 108.1 to
152.6~seconds, respectively. Its longer activation interval and higher
resource cost in Figure~\ref{fig:backlog-mean-costs} therefore matter
alongside the faster recovery. These observations support comparing
latency, recovery, and cost separately when assessing a response.

Kafka scaling had the lowest mean run-level p99 in both conditions,
but its smaller-backlog runs differed substantially in SLO violations
and recovery. Only 1 of those 2 runs met the recovery criterion during
continuing input, so the open recovery marker shows that individual
result and no two-run mean is reported. The run without confirmed recovery also had a more concentrated final
partition distribution than the larger-backlog runs, which both
recovered. Their contrast therefore cannot be attributed to starting
backlog alone.

\begin{figure}[htbp]
\centering
\includegraphics[width=\linewidth,height=.78\textheight,keepaspectratio]{figures/backlog-mean-outcomes.pdf}
\caption{Evaluation-message outcomes and confirmed recovery times.}
\label{fig:backlog-mean-outcomes}
\end{figure}

Consumer CPU and memory costs cover evaluation and drain and include
only observed intervals. Measurement coverage exceeded 99\% of the
recorded consumer operating time in these runs. Redistribution had
lower CPU and memory costs than targeted scaling in both backlog
conditions, while completing all evaluation messages. Keep~3 used
less CPU but completed fewer messages, so its lower CPU cost does
not establish a more efficient response.

\begin{figure}[htbp]
\centering
\includegraphics[width=\linewidth,height=.70\textheight,keepaspectratio]{figures/backlog-mean-costs.pdf}
\caption{Consumer resource cost and decision-to-active time.}
\label{fig:backlog-mean-costs}
\end{figure}
\FloatBarrier
